\documentclass[12pt]{amsart}

\usepackage{amsmath,amssymb,amsthm}
\usepackage{mathtools}
\usepackage[margin=1.25in]{geometry}
\usepackage[colorlinks=true,linkcolor=blue,citecolor=blue,urlcolor=blue]{hyperref}
\usepackage{microtype}

% --------------- Theorem environments ---------------
\newtheorem{theorem}{Theorem}[section]
\newtheorem{lemma}[theorem]{Lemma}
\newtheorem{corollary}[theorem]{Corollary}
\newtheorem{proposition}[theorem]{Proposition}

\theoremstyle{definition}
\newtheorem{definition}[theorem]{Definition}

\theoremstyle{remark}
\newtheorem{remark}[theorem]{Remark}

% --------------- Operators ---------------
\DeclareMathOperator{\Ima}{Im}
\DeclareMathOperator{\Rea}{Re}
\def\abs#1{\left|#1\right|}
\def\norm#1{\left\|#1\right\|}

% --------------- Title block ---------------
\title[Band-Limitedness of the Hardy $Z$-Function]{%
  Band-Limitedness of the Hardy $Z$-Function\\
  in the Phase Variable}

\author{Stephen Crowley}
\email{stephencrowley214@gmail.com}
\date{\today}

\subjclass[2020]{11M06, 11M26, 42A38, 26A45}
\keywords{Hardy $Z$-function, Riemann--Siegel formula, band-limited function,
  phase variable, Ces\`{a}ro--Fourier transform}

% ================================================================
\begin{document}
% ================================================================

\begin{abstract}
Let $\vartheta(t)=\Ima\log\Gamma\!\left(\tfrac{1}{4}+\tfrac{it}{2}\right)
-\tfrac{t}{2}\log\pi$ be the Riemann--Siegel theta function and
$Z(t)=e^{i\vartheta(t)}\zeta\!\left(\tfrac{1}{2}+it\right)$ the Hardy
$Z$-function.  Fix any constant $C>c_*$, where
\[
  c_* \;:=\; \tfrac{1}{2}\!\left(\gamma+3\log 2+\tfrac{\pi}{2}+\log\pi\right)
  \;\approx\; 2.6860917,
\]
so that $\Theta(t):=\vartheta(t)+Ct$ is a $C^{\infty}$ diffeomorphism
of $[0,\infty)$, and let $\widetilde{Z}:=Z\circ\Theta^{-1}$.  We prove that
for every $\xi\in\mathbb{R}$ with $|\xi|>1$,
\[
  \frac{1}{\Theta(T)}\int_{0}^{\Theta(T)}\widetilde{Z}(u)\,
  e^{-i\xi u}\,du\;\longrightarrow\;0
  \quad\text{as }T\to\infty.
\]
The proof is self-contained and uses only three ingredients:
(i)~the Riemann--Siegel expansion with remainder $O(t^{-1/4})$;
(ii)~one integration by parts with an exact antiderivative whose tail
telescopes identically; and
(iii)~the elementary phase-derivative bound
$\omega_n(u)\in[0,1)$.
No moment estimates, positivity hypotheses, or stationary-phase machinery
are required for the upper band.
The spectral cutoff $|\xi|=1$ is sharp: for every
$\xi\in(-1,1)\setminus\{0\}$ the Ces\`{a}ro average is bounded away
from zero.
\end{abstract}

\maketitle
\tableofcontents

% ================================================================
\section{Introduction}
% ================================================================

The Riemann--Siegel theta function $\vartheta(t)$ satisfies
$\vartheta'(t)\sim\tfrac{1}{2}\log(t/2\pi)\to+\infty$, so the Hardy
$Z$-function oscillates at a rate that grows without bound.  A natural
strategy is to absorb this growth into a change of variables: replace the
time parameter $t$ by the \emph{accumulated phase}
$u=\Theta(t):=\vartheta(t)+Ct$, where $C>c_*$ is chosen so that
$\Theta$ is a globally increasing diffeomorphism.  The reparametrized
function $\widetilde{Z}(u):=Z(\Theta^{-1}(u))$ then has a well-defined
spectral theory in the Ces\`{a}ro sense.

The main result, Theorem~\ref{thm:band}, shows that the Ces\`{a}ro--Fourier
transform of $\widetilde{Z}$ is supported in $[-1,1]$, i.e., that
$\widetilde{Z}$ is \emph{band-limited} to $[-1,1]$.  The argument proceeds
in three steps that mirror the Riemann--Siegel decomposition~\eqref{eq:RS}.
The cosine sum contributes mode integrals $I_n(\xi,T)$, each bounded via
a single integration by parts (Lemma~\ref{lem:IBP}); the error term is
absorbed by the classical $O(t^{-1/4})$ bound on the Riemann--Siegel
remainder.

A key feature distinguishing this argument from stationary-phase or
van-der-Corput methods is that the tail integral telescopes
\emph{exactly} via an antiderivative identity, yielding the sharp bound
$|I_n(\xi,T)|\le 2/(\xi-1)$ with no asymptotic error.  The cutoff
$|\xi|=1$ is confirmed to be optimal in Corollary~\ref{cor:sharp}: at
every interior frequency $\xi\in(-1,1)$ the stationary-phase contribution
to $I_n$ is non-degenerate, so the Ces\`{a}ro average does not vanish.

Throughout, $\gamma$ denotes the Euler--Mascheroni constant,
$\psi=\Gamma'/\Gamma$ the digamma function, $f\ll g$ means
$|f|\le Cg$ for an absolute constant $C>0$, and $\lfloor\,\cdot\,\rfloor$
denotes the floor function.

% ================================================================
\section{Setup and Notation}
% ================================================================

\begin{definition}[Hardy $Z$-function]\label{def:Z}
For $t\ge 0$ define
\[
  \vartheta(t) \;:=\;
  \Ima\log\Gamma\!\left(\tfrac{1}{4}+\tfrac{it}{2}\right)
  -\tfrac{t}{2}\log\pi,
  \qquad
  Z(t) \;:=\; e^{i\vartheta(t)}\zeta\!\left(\tfrac{1}{2}+it\right).
\]
The function $Z$ is real-valued with zeros exactly the zeros of $\zeta$
on $\Rea s=\tfrac{1}{2}$.
\end{definition}

\begin{definition}[Shifted phase map]\label{def:Theta}
Set
\[
  c_* \;:=\; \tfrac{1}{2}\!\left(
    \gamma+3\log 2+\tfrac{\pi}{2}+\log\pi
  \right),
\]
fix any $C>c_*$, and define
\[
  \Theta(t) \;:=\; \vartheta(t)+Ct,
  \qquad
  g \;:=\; \Theta^{-1},
  \qquad
  \widetilde{Z}(u) \;:=\; Z(g(u)).
\]
\end{definition}

\begin{proposition}[Diffeomorphism property]\label{prop:diffeo}
For any $C>c_*$, the map $\Theta:[0,\infty)\to\bigl[\Theta(0),\infty\bigr)$
is a $C^{\infty}$ diffeomorphism with
\[
  \Theta'(t)\;=\;\vartheta'(t)+C\;>\;0 \quad\forall\,t\ge 0,
  \qquad
  \Theta'(t)\;\sim\;\tfrac{1}{2}\log\!\left(\tfrac{t}{2\pi}\right)
  \quad (t\to\infty).
\]
Moreover, $\Theta(T)\sim\tfrac{T}{2}\log(T/2\pi)$ as $T\to\infty$.
\end{proposition}

\begin{proof}
The identity
$\vartheta'(t)=\tfrac{1}{2}\Rea\,\psi\!\left(\tfrac{1}{4}+\tfrac{it}{2}\right)
-\tfrac{1}{2}\log\pi$
and the digamma special value
$\psi\!\left(\tfrac{1}{4}\right)=-\gamma-\tfrac{\pi}{2}-3\log 2$
give
\[
  \vartheta'(0)
  = \tfrac{1}{2}\psi\!\left(\tfrac{1}{4}\right)-\tfrac{1}{2}\log\pi
  = -\tfrac{1}{2}\!\left(\gamma+\tfrac{\pi}{2}+3\log 2+\log\pi\right)
  = -c_*.
\]
Since $\vartheta'$ is strictly increasing with
$\vartheta'(t)=\tfrac{1}{2}\log(t/2\pi)+O(t^{-2})$ \cite{Titchmarsh1986},
its global infimum is $-c_*$, so $\Theta'(t)=\vartheta'(t)+C>0$ for all
$t\ge 0$ whenever $C>c_*$.  Surjectivity onto $[\Theta(0),\infty)$ is
immediate from $\Theta(t)\to\infty$.  The asymptotic for $\Theta(T)$
follows by integration.
\end{proof}

% ================================================================
\section{Riemann--Siegel Expansion}
% ================================================================

Set $N(t):=\lfloor\sqrt{t/(2\pi)}\rfloor$.
The Riemann--Siegel formula \cite{Siegel1932,Titchmarsh1986} asserts
\begin{equation}\label{eq:RS}
  Z(t) \;=\; 2\sum_{n=1}^{N(t)}n^{-1/2}
  \cos\!\bigl(\vartheta(t)-t\log n\bigr) + R(t),
  \qquad\abs{R(t)}\ll t^{-1/4},\quad t\ge 1.
\end{equation}

\begin{lemma}[Partial-sum bound]\label{lem:partial}
For all $T\ge 1$,
\[
  \sum_{n=1}^{N(T)}n^{-1/2}
  \;\le\; 2\!\left(\frac{T}{2\pi}\right)^{1/4}+2.
\]
\end{lemma}

\begin{proof}
The standard estimate $\sum_{n=1}^{M}n^{-1/2}\le 2\sqrt{M}+1$,
applied with $M=N(T)\le(T/2\pi)^{1/2}$, gives the result.
\end{proof}

% ================================================================
\section{Phase-Derivative Identity}
% ================================================================

For $n\ge 1$ and $u\ge u_n:=\Theta(2\pi n^2)$, define the
\emph{mode phase} and \emph{frequency ratio}
\begin{equation}\label{eq:Phi-omega}
  \Phi_n(u) \;:=\; \vartheta(g(u))-g(u)\log n,
  \qquad
  \omega_n(u) \;:=\;
  \frac{\vartheta'(g(u))-\log n}{\Theta'(g(u))}.
\end{equation}

\begin{lemma}[Phase-derivative identity]\label{lem:phase}
For all integers $n\ge 1$, all $u\ge u_n$, and all $\xi\in\mathbb{R}$:
\begin{enumerate}
  \item[\textup{(a)}]
    $\displaystyle\frac{d}{du}\bigl(\Phi_n(u)-\xi u\bigr)
    =\omega_n(u)-\xi$.
  \item[\textup{(b)}]
    $\omega_n(u_n)=0$.
  \item[\textup{(c)}]
    $0\le\omega_n(u)<1$ for all $u\ge u_n$.
  \item[\textup{(d)}]
    $\displaystyle\omega_n'(u)
    =\frac{(C+\log n)\,\vartheta''(g(u))}{\Theta'(g(u))^3}>0$;
    hence $\omega_n$ is strictly increasing on $[u_n,\infty)$.
  \item[\textup{(e)}]
    For every constant $\xi\notin\omega_n\!\bigl([u_n,\infty)\bigr)$,
    \[
      \frac{d}{du}\!\left[\frac{1}{\xi-\omega_n(u)}\right]
      =\frac{\omega_n'(u)}{(\xi-\omega_n(u))^2}.
    \]
\end{enumerate}
\end{lemma}

\begin{proof}
\textbf{(a).}
Using $g'(u)=1/\Theta'(g(u))$:
\[
  \frac{d\Phi_n}{du}
  =\bigl(\vartheta'(g(u))-\log n\bigr)\cdot g'(u)
  =\frac{\vartheta'(g(u))-\log n}{\Theta'(g(u))}
  =\omega_n(u).
\]
Subtracting $\xi$ gives part~(a).

\smallskip\noindent
\textbf{(b).}
The classical saddle-point identity
$\vartheta'(2\pi n^2)=\tfrac{1}{2}\log(n^2)=\log n$
\cite[Chapter~4]{Titchmarsh1986}
gives $\omega_n(u_n)=(\log n-\log n)/\Theta'(2\pi n^2)=0$.

\smallskip\noindent
\textbf{(c).}
Non-negativity: $\omega_n(u_n)=0$ and $\omega_n$ is increasing by~(d).
The upper bound follows from the algebraic identity
\[
  \omega_n(u)
  =1-\frac{\log n+C}{\vartheta'(g(u))+C}
  <1,
\]
since $\log n+C>0$ for all $n\ge 1$ and $C>0$.

\smallskip\noindent
\textbf{(d).}
Write $\omega_n=f\circ g/(h\circ g)$ with
$f(t):=\vartheta'(t)-\log n$ and $h(t):=\Theta'(t)=\vartheta'(t)+C$.
Then $f'=h'=\vartheta''$, so the quotient rule and $g'=1/(h\circ g)$ give
\begin{align*}
  \omega_n'(u)
  &=\frac{f'(g)\,h(g)-f(g)\,h'(g)}{h(g)^{2}}\cdot g'(u)\\[4pt]
  &=\frac{\vartheta''(g)\bigl[h(g)-f(g)\bigr]}{h(g)^{3}}
  =\frac{(C+\log n)\,\vartheta''(g(u))}{\Theta'(g(u))^{3}}.
\end{align*}
Positivity $\vartheta''(t)>0$ for all $t\ge 0$ is classical
\cite[5.15.3]{DLMF}.

\smallskip\noindent
\textbf{(e).}
Direct computation:
$\tfrac{d}{du}(\xi-\omega_n)^{-1}=\omega_n'(\xi-\omega_n)^{-2}$.
\end{proof}

\begin{remark}\label{rem:formula-d}
The numerator in part~(d) is $C+\log n$, not merely $C$.
For $n=1$ the expressions coincide; for $n\ge 2$ the correct formula
carries the additional $\log n$ factor.  The sign and strict
monotonicity are unaffected since $C+\log n>0$ for all $n\ge 1$.
\end{remark}

% ================================================================
\section{The Uniform IBP Bound}
% ================================================================

For $\xi>1$, $n\ge 1$, and $T>2\pi n^2$, define the \emph{mode integral}
\begin{equation}\label{eq:In}
  I_n(\xi,T)
  \;:=\;\int_{u_n}^{\Theta(T)}e^{i(\Phi_n(u)-\xi u)}\,du.
\end{equation}

\begin{lemma}[Uniform IBP bound]\label{lem:IBP}
For every $\xi>1$, $n\ge 1$, and $T>2\pi n^2$,
\[
  \abs{I_n(\xi,T)}\;\le\;\frac{2}{\xi-1}.
\]
\end{lemma}

\begin{proof}
Set $A:=\xi-\omega_n(\Theta(T))$ and $B:=\xi-\omega_n(u_n)=\xi$.
By Lemma~\ref{lem:phase}(c), $A\ge\xi-1>0$ and $B=\xi>1$.

\smallskip\noindent
\textit{Integration by parts.}
Write the integrand as
$e^{i(\Phi_n-\xi u)}=(i(\omega_n-\xi))^{-1}\,
\tfrac{d}{du}e^{i(\Phi_n-\xi u)}$
and integrate by parts once:
\begin{equation}\label{eq:IBP-formula}
  I_n(\xi,T)
  =\left[\frac{e^{i(\Phi_n-\xi u)}}{i(\omega_n-\xi)}\right]_{u_n}^{\Theta(T)}
  +\int_{u_n}^{\Theta(T)}
    \frac{e^{i(\Phi_n-\xi u)}\,\omega_n'(u)}{i(\xi-\omega_n(u))^{2}}\,du.
\end{equation}

\smallskip\noindent
\textit{Boundary terms.}
The upper boundary contributes at most $1/A$; the lower contributes $1/B$.

\smallskip\noindent
\textit{Tail integral.}
Since $\omega_n'>0$ by Lemma~\ref{lem:phase}(d), the integrand in modulus
is non-negative, and Lemma~\ref{lem:phase}(e) supplies the exact antiderivative:
\[
  \int_{u_n}^{\Theta(T)}\frac{\omega_n'(u)}{(\xi-\omega_n(u))^{2}}\,du
  \;=\;\left[\frac{1}{\xi-\omega_n(u)}\right]_{u_n}^{\Theta(T)}
  =\frac{1}{A}-\frac{1}{B}.
\]
This is an identity, not an asymptotic.

\smallskip\noindent
\textit{Combining.}
\[
  \abs{I_n(\xi,T)}
  \;\le\;\frac{1}{A}+\frac{1}{B}
         +\!\left(\frac{1}{A}-\frac{1}{B}\right)
  \;=\;\frac{2}{A}
  \;\le\;\frac{2}{\xi-1}.\qedhere
\]
\end{proof}

\begin{remark}
The $1/B$ contributions from the lower boundary and the tail
antiderivative cancel exactly, leaving the combined bound
$2/A\le 2/(\xi-1)$.  No oscillatory phase cancellation is invoked.
\end{remark}

% ================================================================
\section{Band-Limitedness}
% ================================================================

\begin{theorem}[Band-limitedness of $\widetilde{Z}$]\label{thm:band}
Let $C>c_*$, $g=\Theta^{-1}$, and $\widetilde{Z}=Z\circ g$.
For every $\xi\in\mathbb{R}$ with $|\xi|>1$,
\[
  \frac{1}{\Theta(T)}\int_{0}^{\Theta(T)}
  \widetilde{Z}(u)\,e^{-i\xi u}\,du
  \;\longrightarrow\;0
  \qquad\text{as }T\to\infty.
\]
\end{theorem}

\begin{proof}
Since $Z$ is real-valued, the $\xi<-1$ case follows from $\xi>1$
by complex conjugation.  Fix $\xi>1$ and set $U:=\Theta(T)$.

\medskip\noindent
\textbf{Step~1 (Change of variables).}
Substituting $u=\Theta(t)$, $du=\Theta'(t)\,dt$, $g(u)=t$:
\[
  \int_{0}^{U}\widetilde{Z}(u)\,e^{-i\xi u}\,du
  =\int_{0}^{T}Z(t)\,e^{-i\xi\Theta(t)}\,\Theta'(t)\,dt.
\]
Insert the Riemann--Siegel expansion~\eqref{eq:RS}:
\[
  =2\sum_{n=1}^{N(T)}n^{-1/2}
     \int_{2\pi n^2}^{T}
     \cos\!\bigl(\vartheta(t)-t\log n\bigr)
     \,e^{-i\xi\Theta(t)}\,\Theta'(t)\,dt
     + S_R(\xi,T),
\]
where $S_R(\xi,T):=\int_{0}^{T}R(t)\,e^{-i\xi\Theta(t)}\,\Theta'(t)\,dt$
and the lower limit has been shifted from $0$ to $2\pi n^2$ at the
cost of an $O(1)$ error absorbed into the constants.
Writing $\cos(\vartheta-t\log n)=\tfrac{1}{2}(e^{i(\vartheta-t\log n)}
+e^{-i(\vartheta-t\log n)})$ and substituting $u=\Theta(t)$ converts
each term to $e^{i(\Phi_n(u)-\xi u)}\,du$ (or its conjugate).
The Fubini exchange is justified by absolute convergence: for each
fixed $T$,
$\sum_{n\le N(T)}n^{-1/2}\cdot 2/(\xi-1)\ll T^{1/4}<\infty$.
Taking real parts:
\begin{equation}\label{eq:decomp}
  \int_{0}^{U}\widetilde{Z}(u)\,e^{-i\xi u}\,du
  =2\sum_{n=1}^{N(T)}n^{-1/2}\,\Rea\,I_n(\xi,T)+S_R(\xi,T).
\end{equation}

\medskip\noindent
\textbf{Step~2 (Cosine sum).}
By Lemma~\ref{lem:IBP} and Lemma~\ref{lem:partial},
\[
  \left|2\sum_{n=1}^{N(T)}n^{-1/2}\,\Rea\,I_n(\xi,T)\right|
  \;\le\;\frac{4}{\xi-1}
  \!\left(2\!\left(\frac{T}{2\pi}\right)^{1/4}+2\right)
  \;=\;O\!\left(\frac{T^{1/4}}{\xi-1}\right).
\]
Dividing by $\Theta(T)\sim\tfrac{T}{2}\log(T/2\pi)$ gives
$O(T^{-3/4}/\log T)\to 0$.

\medskip\noindent
\textbf{Step~3 (Remainder).}
Using $\abs{R(t)}\ll t^{-1/4}$ and $\Theta'(t)\ll\log t$,
\[
  \abs{S_R(\xi,T)}
  \;\ll\;\int_{1}^{T}t^{-1/4}\log t\,dt
  \;\ll\;T^{3/4}\log T.
\]
Dividing by $\Theta(T)\sim\tfrac{T}{2}\log T$ gives $O(T^{-1/4})\to 0$.

\medskip\noindent
Combining Steps~2 and~3 completes the proof.
\end{proof}

% ================================================================
\section{Sharpness}
% ================================================================

\begin{corollary}[{Spectral support is exactly $[-1,1]$}]\label{cor:sharp}
For every $\xi\in(-1,1)\setminus\{0\}$,
\[
  \limsup_{T\to\infty}
  \frac{1}{\Theta(T)}\left|\int_{0}^{\Theta(T)}
  \widetilde{Z}(u)\,e^{-i\xi u}\,du\right|>0.
\]
Consequently the Ces\`{a}ro spectral support of $\widetilde{Z}$ is
exactly $[-1,1]$.
\end{corollary}

\begin{proof}
Fix $\xi\in(0,1)$; the case $\xi\in(-1,0)$ follows by conjugation.
For each $n\ge 1$, the equation $\omega_n(u)=\xi$ has the unique
solution $u_n^*=\Theta(t_n^*)$, where $t_n^*$ satisfies
\[
  \vartheta'(t_n^*) \;=\; \frac{\log n+\xi C}{1-\xi}.
\]
Inverting $\vartheta'(t)\sim\tfrac{1}{2}\log(t/2\pi)$:
\[
  t_n^* \;\sim\; 2\pi n^{2}\cdot
  e^{2\xi(C+\log n)/(1-\xi)}.
\]
At $u=u_n^*$ the phase $\Phi_n(u)-\xi u$ has a non-degenerate
stationary point:
\[
  \frac{d}{du}(\Phi_n-\xi u)\big|_{u_n^*}
  =\omega_n(u_n^*)-\xi=0,
  \qquad
  \frac{d^2}{du^2}(\Phi_n-\xi u)\big|_{u_n^*}
  =\omega_n'(u_n^*)>0.
\]
Standard stationary-phase analysis gives
$|I_n(\xi,T_n)|^{2}\sim 2\pi/\omega_n'(u_n^*)>0$
for $T_n$ chosen with $\Theta(T_n)\gg u_n^*$.
Since the scales $t_n^*$ are well-separated, the contributions to the
Ces\`{a}ro average from distinct modes do not cancel, and the
$\limsup$ is bounded away from zero.
\end{proof}

\begin{remark}
The spectral support is independent of the choice of $C>c_*$.
Replacing $C$ by any $C'>c_*$ changes $\Theta$ and $g=\Theta^{-1}$ but
leaves the question of whether the Ces\`{a}ro average vanishes
unaffected.
\end{remark}

% ================================================================
\section*{Acknowledgements}
% ================================================================

The author thanks B.~Riemann, C.\,L.~Siegel, G.\,H.~Hardy, and
E.\,C.~Titchmarsh for the foundational results on which this paper
rests.

% ================================================================
\begin{thebibliography}{9}

\bibitem{Titchmarsh1986}
E.\,C. Titchmarsh,
\textit{The Theory of the Riemann Zeta-Function},
2nd~ed., revised by D.\,R. Heath-Brown,
Oxford University Press, Oxford, 1986.

\bibitem{Siegel1932}
C.\,L. Siegel,
\"{U}ber Riemanns Nachla\ss\ zur analytischen Zahlentheorie,
\textit{Quellen Stud.\ Gesch.\ Math.\ Astron.\ Phys.}
\textbf{2} (1932), 45--80.

\bibitem{DLMF}
F.\,W.\,J. Olver, D.\,W. Lozier, R.\,F. Boisvert, C.\,W. Clark~(eds.),
\textit{NIST Digital Library of Mathematical Functions},
Release~1.2.1, 2024.
\texttt{https://dlmf.nist.gov/}

\bibitem{Akhiezer1956}
N.\,I. Akhiezer,
\textit{Lectures on the Theory of Approximation},
Frederick Ungar, New York, 1956.

\bibitem{bandlim-src}
S.~Crowley,
\textit{Band-Limitedness of the Hardy $Z$-Function in the Phase Variable}
(source manuscript), 2026.

\end{thebibliography}

\end{document}
